\section*{Chapitre \ref{ch:g12:synthesis-strategy} --- Stratégie de synthèse et chimie verte}

\begin{solution}{exo:g12:synthesis-strategy:1}
Hydratation~: $46.0/(28.0 + 18.0) = \qty{100}{\percent}$. Fermentation~:
$2 \times 46.0/180.0 = \qty{51.1}{\percent}$.
\end{solution}

\begin{solution}{exo:g12:synthesis-strategy:2}
Les deux groupes, \ce{Z} sur l'\omterm{def:g11:functional-groups:amine}{amine} de la glycine et l'\omterm{def:g11:functional-groups:carboxylic-acid}{ester} méthylique
sur l'acide
de l'alanine, sont mis en place à l'étape 1 et retirés à l'étape 3.
\end{solution}

\begin{solution}{exo:g12:synthesis-strategy:3}
Principe 5~: utiliser des \omterm{def:g6:solutions-and-solubility:solute}{solvants} et des conditions de réaction plus
sûrs.
\end{solution}

\begin{solution}{exo:g12:synthesis-strategy:4}
Un excès d'éthanol rend $Q < K$, si bien que la réaction progresse
davantage dans le sens direct. Un \omterm{def:g12:catalysis:catalyst}{catalyseur} accélère les deux sens de
la même façon~: on atteint le même \omterm{def:g10:reaction-progress-table:system}{état final}, seulement plus tôt.
\end{solution}

\begin{solution}{exo:g12:synthesis-strategy:5}
Oui~: seule la \omterm{def:g11:functional-groups:carbonyl}{cétone} est réduite, en groupe hydroxyle d'\omterm{def:g11:functional-groups:alcohol}{alcool}
secondaire~; l'\omterm{def:g11:functional-groups:carboxylic-acid}{ester} est inchangé.
\end{solution}

\begin{solution}{exo:g12:synthesis-strategy:6}
Plus d'\omterm{def:g11:functional-groups:carboxylic-acid}{ester}~: un excès de l'un des \omterm{def:g7:chemical-reactions:reactant}{réactifs}, ou l'élimination de l'eau
(piège à eau) au fur et à mesure de sa formation. Plus vite~: chauffer à
reflux, ajouter un \omterm{def:g12:catalysis:catalyst}{catalyseur} acide.
\end{solution}

\begin{solution}{exo:g12:synthesis-strategy:7}
Avec \ce{NaBH4}~:
\[ \ce{CH3-CH(OH)-CH2-CH2-CO-O-CH3} . \]
Avec \ce{LiAlH4}~:
\[ \ce{CH3-CH(OH)-CH2-CH2-CH2OH} \]
(pentane-1,4-diol), plus du méthanol provenant de l'autre moitié de
l'\omterm{def:g11:functional-groups:carboxylic-acid}{ester}.
\end{solution}

\begin{solution}{exo:g12:synthesis-strategy:8}
L'\omterm{def:g12:curly-arrows:categories}{addition} (\qty{100}{\percent}), l'estérification
(\qty{83}{\percent}) et la \omterm{def:g10:synthesis-yield:synthesis}{synthèse} de l'ibuprofène en trois étapes
(\qty{77}{\percent}). $77.4/40.0 = 1.9$~: presque deux fois meilleure.
\end{solution}

\begin{solution}{exo:g12:synthesis-strategy:9}
$180.0/(138.0 + 102.0) = \qty{75.0}{\percent}$. Acide éthanoïque~:
$60.0/180.0 = \qty{0.333}{kg}$ par kilogramme d'aspirine.
\end{solution}

\begin{solution}{exo:g12:synthesis-strategy:10}
Le chauffage accélère la réaction~; le réfrigérant renvoie les vapeurs
dans le ballon, si bien qu'aucun \omterm{def:g7:chemical-reactions:reactant}{réactif} ni \omterm{def:g6:solutions-and-solubility:solute}{solvant} n'est perdu. Le
chauffage améliore la vitesse, pas le \omterm{def:g10:synthesis-yield:yield}{rendement} final.
\end{solution}

\begin{solution}{exo:g12:synthesis-strategy:11}
Six étapes contre trois. Principes~: prévenir les déchets (1), maximiser
l'\omterm{def:g12:synthesis-strategy:atom-economy}{économie d'atomes} (2), utiliser des \omterm{def:g12:catalysis:catalyst}{catalyseurs} (9)~; et aussi, moins
d'étapes, moins d'énergie (6).
\end{solution}

\begin{solution}{exo:g12:synthesis-strategy:12}
$0.80 \times 0.70 \times 0.90 = 0.504$~: \qty{50.4}{\percent}. \omterm{def:g7:chemical-reactions:reactant}{Réactif}
de départ~: $0.10/0.504 = \qty{0.20}{mol}$.
\end{solution}

\begin{solution}{exo:g12:synthesis-strategy:13}
Protéger l'\omterm{def:g11:functional-groups:amine}{amine} de l'alanine (\ce{Z}) et l'acide de la glycine (\omterm{def:g11:functional-groups:carboxylic-acid}{ester}
méthylique)~;
former la liaison \omterm{def:g11:functional-groups:amine}{amide} entre l'acide libre de l'alanine et l'\omterm{def:g11:functional-groups:amine}{amine}
libre de la glycine~; retirer les deux \omterm{def:g12:synthesis-strategy:protecting-group}{groupes protecteurs}. Trois
étapes, soit cinq réactions si l'on compte chaque protection et chaque
déprotection.
\end{solution}

\begin{solution}{exo:g12:synthesis-strategy:14}
Le \omterm{def:g12:catalysis:catalyst}{catalyseur} applique le principe 9. Mais \qty{250}{\celsius} va
contre le principe 6, le benzène contre les principes 3 et 5, et une
\omterm{def:g12:synthesis-strategy:atom-economy}{économie d'atomes} de
\qty{35}{\percent} contre le principe 2~: $65/35 = \qty{1.9}{kg}$ de
déchets par kilogramme de produit. Un seul trait vert ne fait pas un
procédé vert.
\end{solution}

\begin{solution}{exo:g12:synthesis-strategy:15}
% ledger: crit:CO2-T, crit:CO2-p
\qty{40}{\celsius} et \qty{100}{bar} sont au-delà du point critique du
dioxyde de carbone, \qty{31}{\celsius} et \qty{73.8}{bar}. La pression
vaut cent fois la pression atmosphérique~: la cuve est en acier épais.
Le dioxyde de carbone n'est ni toxique ni inflammable, il redevient un
gaz sans laisser de résidu, et il peut être réutilisé~; aucun \omterm{def:g6:solutions-and-solubility:solute}{solvant}
chloré ne s'échappe.
\end{solution}

\begin{solution}{pb:g12:synthesis-strategy:1}
% ledger: ibu:bhc-ae-recovered
\textbf{1.} C: $10 + 4 + 1 = 15 = 13 + 2$; H: $14 + 6 + 2 = 22 = 18 + 4$;
O: $3 + 1 = 4 = 2 + 2$.

\textbf{2.} L'acide éthanoïque, \ce{C2H4O2}.

\textbf{3.} L'azote, le chlore et le sodium~: l'ibuprofène ne contient
que du carbone, de l'hydrogène et de l'oxygène.

\textbf{4.} La nouvelle~: principe 9, des \omterm{def:g12:catalysis:catalyst}{catalyseurs} plutôt que des
\omterm{def:g7:chemical-reactions:reactant}{réactifs} stœchiométriques.

\textbf{5.} $13 \times 12.0 + 18 \times 1.0 + 2 \times 16.0 =
\qty{206.0}{g/mol}$.

\textbf{6.} $134.0 + 102.0 + 122.5 + 68.0 + 19.0 + 33.0 + 36.0 =
\qty{514.5}{g/mol}$.

\textbf{7.} $206.0/514.5 = \qty{40.0}{\percent}$.

\textbf{8.} $134.0 + 102.0 + 2.0 + 28.0 = \qty{266.0}{g/mol}$.

\textbf{9.} $206.0/266.0 = \qty{77.4}{\percent}$.

\textbf{10.} $(206.0 + 60.0)/266.0 = \qty{100}{\percent}$ sur le
papier~; plus de
\qty{99}{\percent} en pratique.

\textbf{11.} $514.5/206.0 = \qty{2.50}{t}$.

\textbf{12.} $2.50 - 1 = \qty{1.50}{t}$.

\textbf{13.} $266.0/206.0 = \qty{1.29}{t}$.

\textbf{14.} $60.0/206.0 = \qty{0.29}{t}$.

\textbf{15.} $1.50 - 0.29 = \qty{1.21}{t}$.

\textbf{16.} $0.90^6 = 0.53$: \qty{53}{\percent}.

\textbf{17.} $0.90^3 = 0.73$: \qty{73}{\percent}.

\textbf{18.} Chaque étape utilise des \omterm{def:g6:solutions-and-solubility:solute}{solvants}, des séparations et des
purifications, et perd un peu de produit~: moins d'étapes, moins de ces
pertes.

\textbf{19.} Presque aucune~: son seul sous-produit est utilisé.

\textbf{20.} La voie en trois étapes évite environ \qty{1.5}{t} de
déchets par tonne d'ibuprofène (\qty{1.2}{t} même si son acide
éthanoïque était jeté).
\end{solution}
